\scanpage{60}
\mspage{29}

On \deleted{\unsure{pourra|trouve}} s'intéressera au ss-groupe $G$ de $\unsure{\mathscr G|\mathcal G}$ défini par

\[
 \deleted{\unsure{\mathscr G}}G=\{f\in\mathscr G\mid f\text{ centralise }Z_0
 \inserted{\bmod\unsure{\mathfrak z|z}_0}\ \text{i.e. }\
 \operatorname{int}(f)|Z_0/\mathfrak z_0=\operatorname{id}_{Z_0/\mathfrak z_0}\}
 \tag*{(150)}
\]

Donc on trouve une suite exacte canonique

\diagram{\begin{tikzcd}[ampersand replacement=\&,column sep=1.15em,row sep=1.8em]
{(151)}\&1\arrow[r]\&G_0\arrow[r]\&G\arrow[r]\&Z'_0\arrow[r]\&1\\
\&1\arrow[r,-]\&\mathfrak z_0\arrow[u,arrows={Hooks[harpoon]-},"{\text{cent.}}"']\arrow[r]\&
 Z_0\arrow[u,hook]\arrow[r,-]\& Z_0/\mathfrak z_0\arrow[u,equal]\arrow[r,-]\&1\arrow[u,phantom,"{\unsure{\wr|=}}" description]{,}
\end{tikzcd}}

\deleted{sans que je sois sûr si $G\to Z'_0$ est surjectif} indépendante du choix d'une ``origine'' $A_0\in\Pi$.

Supposons maintenant que le groupe $G_0$ soit transitif sur l'ens \inserted{$A_0$} des arêtes d'une des structures $A_0\in\Pi$, donc de $=$ \inserted{pour} toutes. On voit alors que

\[
 A_0,A_1\in\Pi,\quad A_0\ne A_1\Longrightarrow
 \underbrace{A_0\cap A_1=\emptyset}_{\substack{\text{intersection de}\\
 \text{parties de }\unsure{\mathscr P|\mathcal P|\mathfrak P}_2(S)}}
 \tag*{(152)}
\]

donc \deleted{si} on pose

\[
 A\unsure{:|=} A_\Pi=\bigcup_{\unsure{A'|A_i}\in\Pi}A'\subset\mathscr P_2(S)
 \tag*{(152)}
\]

on a une réunion disjointe donc une application

\[
 A\xrightarrow{\unsure{\tau|t}}\underset{\text{un }Z'_0\text{-torseur}}{\Pi}
 \tag*{(153)}
\]

% archive sheet number: 59
